\documentclass[a4paper,12pt,oneside]{amsbook}
\usepackage[T1]{fontenc}
\usepackage{textcomp,hyperref}
%\usepackage{euler}
%\usepackage{garamond}

\usepackage{newcent}
%\fontfamily{fxm}\selectfont
\usepackage{fancyhdr}
\pagestyle{fancy}
% Detta paket skoter grafiken
\usepackage[dvips]{graphicx}
%\usepackage[pdftex]{graphicx}
\usepackage{tabularx}
\usepackage{amsthm}
\usepackage{amsmath}
\usepackage[swedish]{babel}
% For bortkommentering
\usepackage{verbatim}
% For att slippa indentering av Sections och SubSections
\usepackage{indentfirst}
\usepackage{hyperref}

% Har stalls paragrafindentering och avstand mellan stycken och rader mm
\setlength{\parindent}{0pt}
\setlength{\parskip}{6pt}
\setlength{\baselineskip}{12pt}
% Har definieras sidan
\setlength{\voffset}{-20mm}
\setlength{\hoffset}{-5.5mm}

\setlength{\topmargin}{0mm}
\setlength{\headheight}{10mm}
\setlength{\headsep}{8mm}
\setlength{\footskip}{8mm}

\setlength{\evensidemargin}{0mm}
\setlength{\oddsidemargin}{0mm}

\setlength{\marginparsep}{0mm}
\setlength{\marginparwidth}{0mm}

\setlength{\textwidth}{170mm}
\setlength{\textheight}{240mm}
\setlength{\paperwidth}{210mm}
\setlength{\paperheight}{297mm}
\setlength{\headwidth}{170mm}
% Huvud och Fot
%\fancypagestyle{}{
\fancyhf{}\fancyhead[c]{\small{\textit{Svar och rättningsanvisningar
      till Programmeringsolympiadens skolkval 2025}}}
\fancyfoot[c]{\thepage}
\renewcommand{\headrulewidth}{0.4pt}
\renewcommand{\footrulewidth}{0.4pt}
\newenvironment{lista}
{\begin{itemize}
\setlength{\parindent}{0pt}
\setlength{\itemsep}{6pt}}
{\end{itemize}}
% Uppgifter
\newcounter{probnr}
\newenvironment{tal}{%
\begin{list}
%{\textbf{\arabic{section}.\arabic{probnr}}} {\usecounter{probnr}
{\textbf{\arabic{probnr}}} {\usecounter{probnr}
\setlength{\leftmargin}{0mm}
\setlength{\rightmargin}{0mm}
\setlength{\labelwidth}{-1mm}
\setlength{\labelsep}{1mm}}
\setlength{\itemsep}{6pt}
}{\end{list}}
% Definition av avsnitt: FAKTA, EXEMPEL, PASTAENDE
\newtheorem{fakta}{Fakta}
\newtheorem{exempel}{Exempel}
\newtheorem{problem}{Problem}
%
\newtheoremstyle{test}% NAME
{20pt}      % ABOVESPACE
{10pt}      % BELOWSPACE
{\sffamily} % BODYFONT
{0pt}       % INDENT
{\scshape}  % HEADFONT
{}          % HEADPUNCT
{\newline}  % HEADSPACE
{}          % CUSTOM-HEAD-SPEC

\theoremstyle{test}
\newtheorem{program}{Program}
\newcommand{\sv}[1]{\textsc{#1}}            % Sma versaler
\newcommand{\fe}[1]{\textbf{#1}}            % FET
\newcommand{\ku}[1]{\textit{#1}}            % KURSIV
\newcommand{\cu}[1]{\texttt{#1}}            % Courier
\newcommand{\sk}[1]{\texttt{#1}}            % Courier
\newcommand{\rubrik}[1]{\begin{center}\sf\huge{#1}\normalsize\rm\end{center}}
\begin{document}
%\DeclareGraphicsExtensions{.jpg,.pdf,.mps,.png,.eps}


\specialsection*{Svar och rättningsanvisningar}
\thispagestyle{fancy}
\lhead{}
\begin{itemize}
%\setlength\itemsep{0.2cm}
\item Läs igenom \textit{tävlingsreglerna}.
\item Programmen tas i tur och ordning in i editorn och kompileras.
Uppstår kompileringsfel betraktas programmet som felaktigt och lösningen
ges $0$ poäng.
\item Programmet körs med givna indata enligt nedan. Alternativt, om eleven förberett programmet för det, kan de bifogade indatafilerna användas istället.
\item Varje testfall
  med korrekt svar ger 1 poäng. Svaret anses endast korrekt om det stämmer överens exakt med det anvigna svaret.
\item
Totalt kan man
  alltså få 5 poäng för varje uppgift.
\item Om exekveringstiden för ett testexempel, kört på en modern dator,
\textit{överskrider 3 sekunder} betraktas körningen av testexemplet som felaktigt.
\item Det kan vara viktigt att programmet körs i en miljö liknande den som
programmet utvecklats i, samma version av kompilator eller
interpretator.
\item Vid problem i samband med rättningen är det viktigt att det sunda
förnuftet får råda!
\item Ett förslag till rättningsprocedur kan vara att låta eleven
sitta vid datorn.
\end{itemize}

\vspace{1cm}

\subsection*{Uppgift 1 -- Frack Jack}

.

Det är acceptabelt om deltagaren råkat skriva ut citationstecken i svaret.
Triviala stavfel eller kapitaliseringsfel godtas också.
~\\
{\tt 
\begin{tabular}{||l||c|c||c||}\hline\hline
& \multicolumn{2}{c||}{\fe{Indata}} & \multicolumn{1}{c||}{\fe{Utdata}} \\ 
& Alice kastlängd & Bertils kastlängd &  \\ \hline \hline
\fe{Test 1} & 21 & 3  & Alice\\ \hline
\fe{Test 2} & 10 & 10 & Jack\\ \hline
\fe{Test 3} & 30 & 0  & Bertil\\ \hline
\fe{Test 4} & 22 & 30 & Jack\\ \hline
\fe{Test 5} & 29 & 29 & Jack\\ \hline\hline
\end{tabular}
}
%\vspace{1cm}


\subsection*{Uppgift 2 -- Snurra kompass}
~\\
{\tt 
\begin{tabular}{||l||c|c||c||}\hline\hline
& \multicolumn{2}{c||}{\fe{Indata}} & \multicolumn{1}{c||}{\fe{Utdata}} \\ 
& Gradtal A & Gradtal B &  \\ \hline \hline
\fe{Test 1} & 13  & 256 & 360 \\ \hline
\fe{Test 2} & 180 & 180 & 2 \\ \hline
\fe{Test 3} & 24  & 24  & 15 \\ \hline
\fe{Test 4} & 20  & 34  & 180 \\ \hline
\fe{Test 5} & 147 & 221 & 360 \\ \hline\hline
\end{tabular}
}


\subsection*{Uppgift 3 -- Arborist}

~\\
{\tt 
\begin{tabular}{||l||c|c|l|l||c||}\hline\hline
& \multicolumn{4}{c||}{\fe{Indata}} & \fe{Utdata} \\ 
& N & K & Koordinater & Vikter &  \\ \hline \hline
\fe{Test 1} & 6 & 25 & 10 2 1 25 14 12 & 20 20 20 20 20 20 & 140 \\ \hline
\fe{Test 2} & 7 & 42 & 15 10 29 16 20 20 10 & 20 20 20 20 20 20 20 & 162 \\ \hline
\fe{Test 3} & 5 & 42 & 5 5 5 5 5 & 20 40 20 20 40 & 50 \\ \hline
\fe{Test 4} & 5 & 42 & 29 10 27 25 10 & 20 40 40 20 20 & 162 \\ \hline
\fe{Test 5} & 7 & 42 & 30 29 29 1 10 9 11 & 20 40 20 20 40 20 20 & 176 \\ \hline
\hline\hline
\end{tabular}
}
\vspace{1cm}

\subsection*{Uppgift 4 -- Tingsplatsen}

~\\
{\tt
\begin{tabular}{||l||c|c|l||l||}\hline\hline
& \multicolumn{3}{c||}{\fe{Indata}} & \fe{Utdata} \\
& $N$ & $M$ & Rader & \\ \hline \hline
\fe{Test 1} &1&10&....*...*. & ....*.X.*. \\
\hline
\fe{Test 2} &9&9&......... & ......... \\
&&&.*...*...&.*...*...\\
&&&.........&.........\\
&&&.........&.........\\
&&&.........&...X.....\\
&&&.........&.........\\
&&&.........&.........\\
&&&.*...*...&.*...*...\\
&&&.........&.........\\
\hline
\fe{Test 3} &10&10&.......*.. & .......*.. \\
&&&........*.&........*.\\
&&&..........&..........\\
&&&..........&..........\\
&&&..........&XXXXXXX...\\
&&&..........&.......X..\\
&&&..........&..........\\
&&&..........&..........\\
&&&..........&..........\\
&&&......*...&......*...\\
\hline
\fe{Test 4} &10&9&......... & ....XXXXX \\
&&&.........&....XXXXX\\
&&&.*.......&.*..XXXXX\\
&&&.........&...X.....\\
&&&.........&..X......\\
&&&....*....&XX..*....\\
&&&.........&XX.......\\
&&&.........&XX.......\\
&&&.........&XX.......\\
&&&.........&XX.......\\
\hline
\fe{Test 5} &10&10&..*....... & ..*....... \\
&&&..........&..........\\
&&&*.........&*.........\\
&&&..........&..........\\
&&&......*...&......*...\\
&&&..........&..X.......\\
&&&..........&..........\\
&&&.....*....&.....*....\\
&&&*...*.....&*...*.....\\
&&&.*........&.*........\\
\hline\hline
\end{tabular}
}
\vspace{3cm}

\subsection*{Uppgift 5 -- Långa trappor} .

Precis som trapporna är indatan lång. Detta tvingade oss tyvärr
att dela upp indatan per testfall i två tabeller.

~\\
{\tt 
\begin{tabular}{||l||c|c|c||}\hline\hline
& \multicolumn{3}{c||}{\fe{Indata}} \\ 
& R & C & Bottenraden  \\ \hline \hline
\fe{Test 1} & 923456789 & 2 & 977777777 941111111\\ \hline
\fe{Test 2} & 10000 & 10 & 9999 9999 10004 10014 10008\\
&&& 10004 10002 10004 9999 10002 \\ \hline

\fe{Test 3} & 10000 & 10 & 10004 10003 10011 9999 10005\\
&&& 10001 10004 9999 10004 10006 \\ \hline

\fe{Test 4} & 700000000 & 10 & 699999999 825493196 937342883 998903925 896090514\\
&&& 768426638 852560961 700000006 814808589 700000005 \\ \hline

\fe{Test 5} & 900000000 & 10 & 900000000 922580416 997215756 907792280 947131371\\
&&& 988495167 900000017 900000000 965023170 959759769 \\ \hline

\hline\hline
\end{tabular}
}

\vspace{0.1cm}

Fortsättning...

~\\
{\tt 
\begin{tabular}{||l||c|c|c|c||c||}\hline\hline
& \multicolumn{4}{c||}{\fe{Indata}} & \multicolumn{1}{c||}{\fe{Utdata}} \\ 
& $s_x$ & $s_y$ & $g_x$ & $g_y$ &  \\ \hline \hline
\fe{Test 1} & 1 & 5 & 0 & 911111111 & 442395057470493837 \\ \hline
\fe{Test 2} & 1 & 3218 & 7 & 9996 & 22994427 \\ \hline
\fe{Test 3} & 0 & 7322 & 8 & 9995 & 3592587 \\ \hline
\fe{Test 4} & 2 & 338563662 & 9 & 699999998 & 65318115312042441 \\ \hline
\fe{Test 5} & 1 & 234084934 & 9 & 899999997 & 221721438972343157 \\ \hline
\end{tabular}
}

\vspace{1cm}

\subsection*{Uppgift 6 -- Binära palindrom}

~\\
{\tt 
\begin{tabular}{||l||c|c|c||c||} \hline\hline
  & \multicolumn{3}{c||}{\fe{Indata}} & \fe{Utdata} \\ 
  & l & r & k &  \\ \hline \hline
  \fe{Test 1} & 123456789 & 987654321 & 16 & 6609 \\ \hline
  \fe{Test 2} & 268451841 & 4503599694479361 & 3 & 13 \\ \hline
  \fe{Test 3} & 123456789123456789 & 989999999999999992 & 29 & 42468528 \\ \hline
  \fe{Test 4} & 664545914252227017 & 909090909090909090 & 34 & 28061113 \\ \hline
  \fe{Test 5} & 77777777777 & 510407950671893831 & 31 & 68774807 \\ \hline

\end{tabular}  
}
\vspace{1cm}


\end{document}




\end{document}


